\providecommand{\devnum}[1]{#1}

\section{Additional Theory: Scope, Accounting and Instance-Dependent Bounds}
\label{sec:supp_theory_additions}

The results below delimit the main-text claims rather than enlarge them. The first
example rules out instance optimality for the uniform guard.
The exact identity that follows is a useful job-level accounting decomposition;
its core is standard sample-path workload balance, and we do not treat it as a
new conservation law. The final bounds replace the global timeout by realised
order statistics, but in their actual-size form they are retrospective
certificates under completion-only observation.

\subsection{Comparison with the Nearest Guarantee Disciplines}

\begin{table*}[htbp]
\centering
\caption{Scope of the nearest per-job guarantees and wrapper precedent.}
\label{tab:supp_related_comparison}
\scriptsize
\setlength{\tabcolsep}{2pt}
\begin{tabular}{p{0.11\textwidth}p{0.13\textwidth}p{0.10\textwidth}p{0.08\textwidth}p{0.15\textwidth}p{0.12\textwidth}p{0.18\textwidth}}
\hline
Method & Reference discipline & Arrival assumptions & Preemption & Information needed & Server model & Guarantee \\
\hline
This work & Same-input FCFS & Arbitrary finite path & No & Ranks, completions, public $L$, opaque base proposal & $k$ identical servers & Deterministic per-job additive waiting excess; wrapper around any base \\
PFCFS \cite{schwiegelshohn2000fairness} & Flow time without later submissions & Online releases & Gang preemption & Processor demand; service time need not be known; preemption penalty & Rigid parallel jobs & $\lambda$-fairness; PFCFS is $(2+2\bar s)$-fair with competitive efficiency bounds \\
Conservative backfilling \cite{lindsay2013backfilling} & Submission-time reservation & Online submissions & No & Requested wall time and resources & Rigid parallel jobs & Later backfills do not delay any reserved start \\
PV-EASY \cite{yuan2010pveasy,yuan2014strictfairness} & Priority reservation & Online submissions; prediction error allowed & Kill and restart & Runtime estimates and resources & Rigid parallel jobs & Zero lower-priority delay to the protected reservation \\
Nudge \cite{grosof2021nudge} & FCFS & Stationary Poisson & No & Size class at arrival & One server & Response-time stochastic dominance for its model; at most one swap per job \\
Load Balancing Guardrails \cite{grosof2019guardrails} & Unconstrained dispatcher & Stochastic heavy-traffic regime & SRPT at servers & Size class and dispatch state & Parallel server pools & Any dispatcher becomes mean-response-time optimal in heavy traffic \\
\hline
\end{tabular}
\end{table*}
% Source for the PFCFS definition and constants, and reading-depth notes for all
% rows: prechecks/guard_optimality_verify/literature_check.md and
% prechecks/guard_optimality_verify/verification.md, Section 9.

\subsection{No Uniform Instance-Optimality Ratio}

\begin{proposition}[An unbounded instance-optimality ratio]
\label{prop:supp_unbounded_instance}
At $k=1$, let three batch jobs have rank-order service times $(L,1,1)$, with
$L\geq2$, and set $G=L$. The cap of the parameter rule of the manuscript
(Corollary~3 there, $B_{\max} = k(G-(3-2/k)L)$) is $B=0$, so the
guard is FCFS and has total waiting time $2L+1$. The order $(1,1,L)$ has total
waiting time $3$ and satisfies the per-job promise $G$. Consequently the ratio
between the guard's total wait and that of a feasible policy is at least
$(2L+1)/3$, which is unbounded as $L\to\infty$.
\end{proposition}
% Source and direct checks at L=10,100,1000,10000:
% prechecks/guard_optimality_verify/out_t1.txt; derivation and scope in
% prechecks/guard_optimality_verify/verification.md, Section 4.

\begin{proof}
FCFS produces waits $0,L,L+1$. Under the competing order the two short jobs wait
$0$ and $1$, and the long job waits $2$. The long job's excess is therefore $2$
and both short jobs have nonpositive excess, so $L\geq2$ makes the order
$G$-feasible. The comparison allows a clairvoyant feasible policy; restricting
the comparator to completion-only policies would be a different question.
\end{proof}

The pointwise-maximal budget $B^{*}=G-L+1$ of Proposition~\ref{thm:k1_maximal}
admits one short job here, giving total wait $L+2$ and ratio $(L+2)/3$; the
conclusion is unchanged.

\subsection{Idle-Capacity Accounting}

Let $t_0$ be the first arrival time, let $b_Q(u)$ be the number of busy servers
under policy $Q$ at time $u$, and define the cumulative unused capacity
\[
 \Gamma_Q(t)=\int_{t_0}^{t}\bigl(k-b_Q(u)\bigr)\,du.
\]
This includes idle capacity caused by having fewer than $k$ jobs in the system;
work conservation excludes only deliberate idling while a job waits.

\begin{lemma}[Exact accounting with idle-capacity correction]
\label{lem:supp_gamma_identity}
For every finite input, every non-preemptive work-conserving policy $P$, and
every job $i$,
\[
 k\bigl(W_P[i]-W_{\mathrm{FCFS}}[i]\bigr)
 =\mathrm{In}_i-\mathrm{Out}_i
  +\Gamma_P(a_i)-\Gamma_{\mathrm{FCFS}}(a_i)
  -\rho_i^P+\rho_i^{\mathrm{FCFS}}.
\]
The identity was checked on \devnum{214{,}735} job-instances, including tied
events, non-batch arrivals, and same-phase overtakers.
\end{lemma}
% Source for the 214,735 checks and their decomposition:
% prechecks/guard_optimality_verify/out_t2.txt and
% prechecks/guard_optimality_verify/verification.md, Sections 7.1--7.2.

\begin{proof}
Let $A(t)$ be cumulative arrived work and $E_Q(t)$ cumulative executed work.
Then
\[
 U_Q(t)=A(t)-E_Q(t)
       =A(t)-k(t-t_0)+\Gamma_Q(t).
\]
At $a_i$, after simultaneous arrivals and before that phase's dispatches, work
from jobs not ranked below $i$ is identical in the two runs. Hence
\[
 R_i^P-R_i^{\mathrm{FCFS}}
 =U_P(a_i)-U_{\mathrm{FCFS}}(a_i)
 =\Gamma_P(a_i)-\Gamma_{\mathrm{FCFS}}(a_i).
\]
Substitution into the two busy-period equalities of the main text (the displayed equality $k\,W_P[i] = R^P_i + \mathrm{In}_i - \mathrm{Out}_i - \rho^P_i$ in the proof of the identity), one for $P$
and one for FCFS, proves the formula. A same-phase overtaker enters both
$\mathrm{In}_i$ and $\rho_i^P$ with its full size and therefore cancels.
\end{proof}

Writing $D_i=\Gamma_P(a_i)-\Gamma_{\mathrm{FCFS}}(a_i)$, the arrival-frozen
budget
\[
 b_i=\max\{0,\ k(G-L)-(k-1)L-D_i\}
\]
has the mathematical promise $G$ in the standard range
$G\geq(3-2/k)L$: combine
$\mathrm{In}_i<b_i+kL$ with Lemma~\ref{lem:supp_gamma_identity} and
$\rho_i^{\mathrm{FCFS}}\leq(k-1)L$. This is not a deployable rule under the
paper's completion-only observation model. In general
$\Gamma_{\mathrm{FCFS}}(a_i)$ depends on when an FCFS-shadow job completes
before its size has been revealed in the realised run; simulating the shadow
does not supply that missing completion time.

\subsection{Instance-Dependent Order-Statistic Bounds}

For a set of jobs $S$, let $\Lambda_r(S)$ be the sum of its largest at most $r$
original service times, with missing terms replaced by zero. We abbreviate
\[
 \Lambda_{k-1}(t)=\Lambda_{k-1}(\{j:a_j\leq t\}),\quad
 \Lambda^{<i}_{k-1}=\Lambda_{k-1}(\{j:j<i\}),
\]
and
\[
 \Lambda^{>i}_{k}(t)=\Lambda_k(\{j:j>i,\ a_j\leq t\}).
\]

\begin{proposition}[Instance-dependent workload gap]
\label{prop:supp_lambda_gap}
For any two non-preemptive work-conserving policies $Q_1,Q_2$,
\[
 |U_{Q_1}(t)-U_{Q_2}(t)|\leq\Lambda_{k-1}(t)
 \qquad\text{for every }t.
\]
\end{proposition}

\begin{proof}
The workload difference is continuous across arrivals, whereas
$\Lambda_{k-1}(t)$ can only increase. If
$U_{Q_1}(t)>\Lambda_{k-1}(t)$, policy $Q_1$ must have at least $k$ jobs present:
otherwise their remaining work is no larger than the sum of their at most
$k-1$ original sizes. All its servers are therefore busy, so the difference
cannot increase. It cannot cross the nondecreasing barrier from below; exchange
$Q_1$ and $Q_2$ for the other direction.
\end{proof}

\begin{proposition}[Instance-dependent excess accounting]
\label{prop:supp_lambda_excess}
For every policy $P$ and job $i$,
\[
 k\,\mathrm{excess}_P[i]
 \leq \mathrm{In}_i-\mathrm{Out}_i
 +\Lambda_{k-1}(a_i)+\Lambda^{<i}_{k-1}.
\]
\end{proposition}

\begin{proof}
In Lemma~\ref{lem:supp_gamma_identity}, Proposition~\ref{prop:supp_lambda_gap}
bounds the idle-capacity difference, equivalently the workload difference at
$a_i$, by $\Lambda_{k-1}(a_i)$. Drop $-\rho_i^P\leq0$. The at most $k-1$ FCFS
jobs still running when $i$ starts all have rank below $i$, so
$\rho_i^{\mathrm{FCFS}}\leq\Lambda^{<i}_{k-1}$.
\end{proof}

\begin{proposition}[Instance-dependent overtaking under the guard]
\label{prop:supp_lambda_in}
If job $i$ has an overtaker and $\tau$ is the dispatch time of its last
overtaker, then
\[
 \mathrm{In}_i
 <\mathrm{budget}(i,\tau)+\Lambda^{>i}_k(s_i^P).
\]
If the budget is nondecreasing in time, the first term may be replaced by
$\mathrm{budget}(i,s_i^P)$; with only a uniform cap it may be replaced by
$B_{\max}$.
\end{proposition}

\begin{proof}
Completed overtaking work before the last overtaker is dispatched is strictly
below $\mathrm{budget}(i,\tau)$. At $\tau$, at most $k$ dispatched overtakers
remain unfinished: the newly dispatched job and at most one on each other
server. Their total original size is at most
$\Lambda^{>i}_k(s_i^P)$.
\end{proof}

\begin{proposition}[Combined instance-dependent guard bound]
\label{prop:supp_lambda_guard}
Under a cap $B_{\max}$, every job satisfies
\[
 k\,\mathrm{excess}_P[i]
 \leq B_{\max}+\Lambda^{>i}_k(s_i^P)
 +\Lambda_{k-1}(a_i)+\Lambda^{<i}_{k-1},
\]
and the inequality is strict when $i$ has an overtaker. When it has none, the
sharper weak bound omits the first two terms.
\end{proposition}

\begin{proof}
Combine Propositions~\ref{prop:supp_lambda_excess} and
\ref{prop:supp_lambda_in}, use $\mathrm{Out}_i\geq0$, and apply the cap. If no
job overtakes $i$, then $\mathrm{In}_i=0$ and
Proposition~\ref{prop:supp_lambda_excess} gives the stated sharper form.
\end{proof}

These $\Lambda$ quantities use actual sizes, including sizes of unfinished
jobs, and hence are offline instance certificates in the completion-only model.
Replacing each unknown size by a public class-specific upper bound gives a
deployable but generally looser inequality. The homogeneous substitution recovers
the main text's aggregate worst-case constant; it does not prove that each
order-statistic term is separately minimal.

As an illustration, consider independently sampled truncated-Pareto sizes with
$\alpha=1.5$, $x_{\min}=1$, hard cap $1000$, sample size $N=200$, and
\devnum{5{,}000} samples drawn with seed \devnum{73621}. The first ratio below
sets $L$ after the fact to the maximum of each finite sample; the second uses the
fixed hard cap as the denominator.
% Source for all parameters and values in this paragraph and table:
% prechecks/guard_optimality_verify/out_t3.txt and
% prechecks/guard_optimality_verify/verification.md, Section 8.3.
\begin{center}
\begin{tabular}{rcc}
\hline
$k$ & $\mathbb{E}[\Lambda_{k-1}/((k-1)L_{\mathrm{sample}})]$
    & $\mathbb{E}[\Lambda_{k-1}/((k-1)L_{\mathrm{cap}})]$ \\
\hline
$4$  & $0.6825$ & $0.0444$ \\
$8$  & $0.4648$ & $0.0263$ \\
$16$ & $0.3139$ & $0.0161$ \\
\hline
\end{tabular}
\end{center}
Thus the larger ratios use $L=L_{\mathrm{sample}}=\max_j C_j$, not the published
timeout $L_{\mathrm{cap}}=1000$.

\subsection{Why Kill-and-Restart Tiering Does Not Escape the Floor}
\label{sec:supp_tiering}

Section~9 of the manuscript states that the obvious escape from the $\Omega(L)$ floor
fails. Kill-and-restart tiering runs every job for a short first quantum $\tau$ and
restarts the survivors, discarding the killed attempt; on $n$ jobs of length $L$ arriving
together it leaves some job with an excess of $n\tau/k - L$ over first-come first-served,
and that excess is wasted work rather than overtaking, so no budget removes it. The
light-job constant does not improve either, a second-tier piece being still a
non-preemptive block of length $L$: holding $\tau$ fixed and varying $L$, the worst
light-job additive constant we measure grows as about \devnum{1.1}$L$ and is flat in
$\tau$. The mechanism is TAGS \cite{harcholbalter2002tags} with the dedicated hosts
removed, and removing them is what breaks it: in TAGS the kill is a routing device and
the work is redone on the next host, whereas in a judging queue the kill is terminal.
% Ninth pass: this subsection stood in Section 9 of the manuscript, which keeps the
%   statement, the citation and the reason the mechanism breaks.

\subsection{Reservation-and-Refund Charging under Enforced Run-Time Caps}
\label{sec:supp_reservation}

Everything below assumes an information model the manuscript does not. Each job $j$
carries a cap $\ell_j$ with $C_j \le \ell_j \le L$ that is known when the job arrives and
that the service enforces, whereas Sections~5 and~6 of the manuscript assume a service
that learns a job's cost only at its completion. That is why this is a supplementary
subsection and not a
theorem of Section~6: admitting it into the main line would put two information models
side by side.

The mechanism itself is older than the guard, and we claim none of it. Conservative
backfilling gives every queued job a reservation when it is submitted and lets a job move
ahead only if it does not delay any job in the queue \cite{mualem2001backfilling}, and
enforced caps are standard there, a job that exceeds its declared run time being killed.
With one server-slot per job the head's reservation is whichever server is free at that
instant, so in our model both conservative and aggressive (EASY) backfilling reduce to
first-come first-served and neither is a comparator here. Lindsay
et~al.~\cite{lindsay2013backfilling} keep the same tentative profile, reserve on the
enforced estimate, compress the profile when a job ends early, and guarantee every job an
absolute start time fixed at its arrival. Reservation with refund is therefore their
mechanism; what differs here is the quantity guaranteed, a per-job bound on the delay
relative to first-come first-served, proved pathwise, with a maximal permission rule at
one server, and the exact constants at $k=1$. Load-balancing guardrails
\cite{grosof2019guardrails} charge dispatched work per size rank and per server and
balance those counters across servers, with the size known at dispatch, no reference
policy and no budget parameter; their Definition~2.1 does adjust a counter, lowering it
to the minimum across servers when a server empties, so it is not a rule without refund,
only one that never refunds at an individual job's completion.
% Source: prechecks/reservation_guard_verify/verification.md, Section 4 (the three
% readings, the corrected pointer to Definition 2.1, the counter reset, and the
% extension of the collapse argument to EASY) and its Section 5 list of statements that
% survive the check.

Write $\mathrm{over}_R[q](t)$ for the charge a waiting job $q$ carries: for every job of
rank above $q$ already dispatched, its cap while it is in service and its true service
time once it has completed.

\begin{proposition}[Reservation charging bounds the overtaken work]
\label{prop:supp_reservation_in}
Accept the base policy's proposal $j$ only when $j$ is the queue head or when
$\mathrm{over}_R[q]+\ell_j \le Z_q$ for every waiting $q$ of rank below $j$, and dispatch
the head otherwise. Then $\mathrm{In}_i \le Z_i$ for every input, every base policy and
every job $i$.
\end{proposition}

\begin{proof}
Rank is decided by $(a_j,\text{index})$, so every job counted in $\mathrm{In}_i$ is
dispatched after $i$ arrives and before $i$ starts, at an epoch where $i$ is waiting and
is not the head; the head exemption therefore does not apply and the test is run with
$q=i$. The charge carried for an already dispatched job is its cap while it runs and its
true service time afterwards, and both are at least its service time, so
$\mathrm{over}_R[i]$ at the dispatch of the $s$-th overtaker is at least
$\sum_{u<s}C_{j_u}$; the test then gives
$\sum_{u\le s}C_{j_u} \le \mathrm{over}_R[i]+\ell_{j_s} \le Z_i$. A rejected epoch
dispatches the head, whose rank is not above $i$, so it never enters $\mathrm{In}_i$. The
third step is where the cap has to be enforced: a job that may exceed $\ell_j$ breaks it
at once.
\end{proof}

\begin{proposition}[The promise under a constant allowance]
\label{prop:supp_reservation_promise}
With $Z \equiv B$, the rule gives $\mathrm{excess} \le B$ at $k=1$, attained, and
$k\,\mathrm{excess} < B+(2-2/k)L$ for $k \ge 2$, strictly. Setting
$B = k\bigl(G-(2-2/k)L\bigr)$ therefore offers the promise $G$.
\end{proposition}

The additive constant of the promise falls from $(3-2/k)L$ to $(2-2/k)L$, a drop of
exactly $L$ at every $k$: the $kL$ overshoot that completed-work charging pays in the
guard's overtaken-work bound disappears, while the $2(k-1)L$ slack of the identity
remains. That residual is still approached and not attained. On the witness family of
Section~S8.2 run at $B=f+L+1$, the ratio $(k\,\mathrm{excess}-B)/L$ is exactly
$2-2^{1-m}$ after $m$ rounds at $k=2$, reaching \devnum{255/128}, and reaches
\devnum{907/243} at $k=3$ and \devnum{323/64} at $k=4$, against the limits $2k-2$ of
\devnum{2}, \devnum{4} and \devnum{6}; a closed-form witness that attains the constant is
not known.
% Source: prechecks/reservation_guard_verify/verification.md, decision table rows C1-c,
% C1-d and C1-e, and out_c1b_indep.txt, which runs the S8.2 witness to m = 8.

Two consequences are worth separating, because they are easy to confuse. The
\emph{entry threshold} does not drop by $L$. An admitted overtaker always has
$\ell_j \le Z_q$, so at $k=1$ an overtaker of cap $L$ becomes admissible as soon as
$G \ge L$, which meets the $L$-scale floor of Proposition~1 of the manuscript with
equality, and for general $k$ the condition is $G \ge (2-1/k)L$ against
$G > (3-2/k)L$ for Algorithm~1 of the manuscript, a difference of $(1-1/k)L$, which is
zero at one server. It is the additive constant of the promise, not the threshold, that
falls by exactly $L$.

At one server the reservation is inert. A dispatch epoch there has an idle server, so no
job of rank above the head is in service and $\mathrm{over}_R$ coincides with the
completed-work charge: over \devnum{255{,}861} readings at $k=1$ the two never differ,
against \devnum{8{,}277} differences in \devnum{110{,}503} readings at $k \ge 2$. The
single-server gain comes entirely from the forward-looking $\ell_j$ term in the admission
test. For a constant allowance only the queue head has to be tested, because
$\mathrm{over}_R[q]$ is non-increasing in rank, so the implementation cost of
Section~5 of the manuscript is unchanged.
% Source: prechecks/reservation_guard_verify/out_overR_eq_overC_k1.txt (reading counts
% and mismatches at k = 1 and k >= 2) and out_c1b_indep.txt part 0 (19,381 head-only
% decisions, no divergence).

\begin{proposition}[One server: closed form and maximal permission]
\label{prop:supp_reservation_k1}
On the integer two-size batch family at one server, around an exact shortest-job-first
proposal sequence, the rule closes
$\min\bigl(1,\ \max(0,\lfloor (G-\ell)/\sigma \rfloor+1)/n\bigr)$ of the FCFS-to-SJF
total-wait gap. This is $\lfloor G/\sigma \rfloor/n$ at $\ell = C$, the size-aware
optimum, and $(\lfloor (G-L)/\sigma \rfloor+1)/n$ at $\ell = L$, the size-oblivious
frontier, so at $\ell=L$ it matches, and does not beat, the sharp completed-work guard.
Moreover, at any reachable history with queue head $h$, a wrapper that reads the caps and
promises $G$ must reject a later-ranked proposal $j$ whenever
$\mathrm{over}_R[h]+\ell_j > G$; the largest constant allowance with promise $G$ is
therefore exactly $B^{*}=G$, against $B^{*}=G-L+1$ for completed-work charging, and the
statement needs no integer lattice.
\end{proposition}

The maximality argument runs through the relabelling of $j$ to its cap. Three points
decide whether it closes. Only $j$ itself and the unfinished dispatched jobs of rank
above $h$ may be relabelled, and at $k=1$ there are none of the latter, so the general
form of the step is not available and is not needed. $W_{\mathrm{FCFS}}[h]$ is unchanged
because $h$ is the head, every job of rank below $h$ has already been dispatched, and
first-come first-served starts $h$ at a time fixed by the arrivals and service times of
jobs ranked below it alone. And the promise has to be quantified over the base policy's
proposal transcript: a base that reads true sizes, exact shortest-job-first among them,
could propose differently after the relabelling, and without that quantifier the
indistinguishability step fails. Nothing is needed about later steps, since $h$ is the
head, $\mathrm{Out}_h$ is always zero and $\mathrm{In}_h$ is non-decreasing.

At $k \ge 2$ the rule is not pointwise maximal, and this is settled rather than open.
Take two servers, four unit jobs arriving together and $G=0$. At the first epoch the rule
refuses every later-ranked proposal, because $\mathrm{over}_R[0]+\ell_1 = 1 > 0$, yet
dispatching job~1 is safe: the second server takes the head in the same instant, the
remaining two start at $t=1$ exactly as under first-come first-served, and every job has
excess $0$ under any continuation. The scale of the gap is not marginal. Of
\devnum{452{,}810} refused actions at $k=2$, \devnum{416{,}368} leave the victim inside
the promise after relabelling, and \devnum{481{,}964} of \devnum{491{,}340} do so at
$k=3$. This is what one expects once $\mathrm{In}$ and excess stop corresponding one to
one on several servers.
% Source: prechecks/reservation_guard_verify/verification.md, row C2-c and Section 2,
% counterexample 1; out_c2b_indep.txt, refused-action and safe-after-relabelling counts
% at k = 2 and k = 3.

Finally, the residual the reservation cannot remove is not an artefact of charging. There
are inputs on which a job has $\mathrm{In}_i=\mathrm{Out}_i=0$ and excess
$L\bigl(1-((k-1)/k)^m\bigr)$, so no charging rule of any kind removes an $\Omega(L)$
term. Its supremum is $L$, approached and not attained, which is the whole of
$(2-2/k)L$ at $k=2$ and only the part of it up to $L$ for $k \ge 3$.
% Source: prechecks/reservation_guard_verify/verification.md, rows C4a-a to C4a-c and
% out_c4b_indep.txt, k = 2 and k = 3 at m = 1..6 with excess/L = 1/2, 3/4, ..., 63/64
% and 1/3, 5/9, ..., 665/729.

The claims above are the ones an independent check re-derived from the written rule and
re-ran against its own implementation: \devnum{1{,}983{,}467} exhaustive runs and
\devnum{10{,}230{,}898} per-job checks give
$\max(\mathrm{In}-Z)=\devnum{0}$ with no violation, the instance form
$k\,\mathrm{excess} \le B+\Lambda_{k-1}(a_i)+\Lambda^{<i}_{k-1}$ holds with zero margin
over \devnum{1{,}488{,}984} checks, the closed form of
Proposition~\ref{prop:supp_reservation_k1} matches on all \devnum{848} family rows
including the corner cases $\ell = C$ and $\ell = L$, and the $k=1$ maximality step
survives \devnum{573{,}252} relabelling tests. Two statements of the earlier study did
not survive that check and are not made here: that the entry threshold drops by $L$, and
that maximality at $k \ge 2$ is open. A third needs a qualifier: the largest excess
attainable at budget $B$ equals $B$ when the cap is tight, $\ell = C$, and is $0$ while
$B < L$ when $\ell = L$, so in general only $\le B$ can be claimed.
% Source: prechecks/reservation_guard_verify/verification.md, decision table rows C1-a,
% C1-f, C2-a, C2-b, C3-d, C2-c and C6-c, with out_c1_indep.txt, out_t32_indep.txt,
% out_c2a_indep.txt and out_c2b_indep.txt.
